
\documentclass[11pt]{article}
\usepackage[a4paper,margin=1in]{geometry}
\usepackage{amsmath,amssymb,amsthm,mathtools}
\usepackage{booktabs}
\usepackage{enumitem}
\usepackage{xcolor}
\usepackage{hyperref}
\hypersetup{colorlinks=true,linkcolor=blue,citecolor=blue,urlcolor=blue}

\newtheorem{theorem}{Theorem}
\newtheorem{lemma}{Lemma}
\newtheorem{proposition}{Proposition}
\newtheorem{definition}{Definition}
\newtheorem{remark}{Remark}

\newcommand{\eps}{\varepsilon}
\newcommand{\XiBD}{\Xi^{\rm reg}_{BD,N}}
\newcommand{\XiBC}{\Xi^{\rm BC}}
\newcommand{\cH}{c_{{\rm hyb},N}}
\newcommand{\gN}{\gamma_N}
\newcommand{\LamN}{\Lambda_{R,N}}

\title{Step 123: Balanced-Length Theorem for the Regularized Müntz/co-Poisson Shadow}
\author{Riemann membrane construction notes}
\date{}

\begin{document}
\maketitle

\section{Purpose}
Steps 119--122 made the Burnol-to-Dirichlet readability problem lawful.  The naive critical-line partial sum of \(\zeta\) was rejected.  The accepted shadow is a regularized Müntz/co-Poisson shadow
\[
Z_X^\omega(s)
=\sum_{n\ge 1}\omega(n/X)n^{-s}-X^{1-s}\Omega(1-s),
\]
with the pole/Müntz correction included.  The associated readability defect is
\[
\delta^{\rm reg}_{BD,N}
=\bigl\|(I-P^{\rm reg}_{D,N})B_NG_{B,N}^{-1/2}\bigr\|.
\]
Step 123 isolates the new quantitative compatibility gate:
\[
\text{the shadow length must be large enough for Müntz accuracy,}
\]
while
\[
\text{the same length must remain short enough for the source moment theorem.}
\]
This is the balanced-length problem.

\section{Abstract data}
Let \(Y_{R,N}\) be a finite window of the residual sector of \(\Xi^{\rm BC}\).  Let \(B_N\) be the Burnol/co-Poisson response synthesis and let \(D_N^{\rm reg}(X_N)\) be the regularized Dirichlet-readable synthesis at length \(X_N\).  The full source route has the form
\[
G_{\mathcal X,N}\succeq \gamma_N H_N,
\qquad
R_N^*H_NR_N\succeq c_{{\rm hyb},N}G_{R,N},
\]
so that
\[
R_N^*G_{\mathcal X,N}R_N\succeq \gamma_Nc_{{\rm hyb},N}G_{R,N}.
\]
The goal is
\[
\gamma_Nc_{{\rm hyb},N}\to\infty
\]
with fixed/exhaustive tail promotion.

\section{Müntz shadow error model}
\begin{definition}[Müntz shadow parameters]
For a Burnol window \(N\), let
\[
T_N = \text{vertical/spectral height},
\qquad
A_N = \text{uniform Burnol atom complexity},
\]
and suppose the contour residual obeys
\[
\tau_{M,N}(X_N)
\le C_M A_N(1+T_N)^\alpha X_N^{-\eta}
\]
for fixed \(C_M,\alpha,\eta>0\).  Let
\[
\epsilon_{B,N}=\text{Burnol geometric exhaustion defect},
\]
\[
\epsilon_{\rm Mell,N}=\text{Mellin/Dirichlet modeling defect},
\]
\[
\kappa_{\rm tail,N}=\text{coefficient/support/tail defect}.
\]
Define
\[
E_N(X_N):=
\epsilon_{B,N}+\epsilon_{\rm Mell,N}+C_MA_N(1+T_N)^\alpha X_N^{-\eta}+\kappa_{\rm tail,N}.
\]
\end{definition}

By Step 122,
\[
\delta^{\rm reg}_{BD,N}\le \epsilon_{\rm Mell,N}+\tau_{M,N}(X_N)+\kappa_{\rm tail,N},
\]
and hence
\[
c_{{\rm hyb},N}\ge 1-E_N(X_N)^2.
\]

\section{Source length cap}
The source moment theorem is not available at arbitrary Dirichlet length.  We encode its admissible range by a scale \(Q_N\) and exponent \(\theta>0\):
\[
X_N\le Q_N^\theta.
\]
For example, in the Heap--Soundararajan lower-bound method, the source polynomials are deliberately short.  In their zeta setup the constructed Dirichlet polynomial has length bounded by a small power of the ambient scale, schematically \(\le T^{k/9}\).  In this note the analogous input is the black-box length cap \(X_N\le Q_N^\theta\).

\section{Balanced-length theorem}
\begin{theorem}[Balanced-length theorem]
Assume:
\begin{enumerate}[label=(\roman*)]
\item the regularized Müntz/co-Poisson shadow satisfies
\[
\delta^{\rm reg}_{BD,N}
\le
\epsilon_{\rm Mell,N}+C_MA_N(1+T_N)^\alpha X_N^{-\eta}+\kappa_{\rm tail,N};
\]
\item the source theorem is valid for lengths \(X_N\le Q_N^\theta\), with
\[
G_{\mathcal X,N}\succeq \gamma_NH_N;
\]
\item the finite residual windows have fixed/exhaustive tail promotion.
\end{enumerate}
If there exists a choice \(X_N\le Q_N^\theta\) such that
\[
E_N(X_N)<1
\]
and
\[
\gamma_N\bigl(1-E_N(X_N)^2\bigr)\to\infty,
\]
then the residual sector is source-absorbed:
\[
\Xi^{\rm BC}\preceq F_N+E_N^{\rm abs}
\]
with vanishing completed defect after tail promotion.
\end{theorem}

\begin{proof}
By the Step 122 angular-gap estimate,
\[
c_{{\rm hyb},N}\ge 1-E_N(X_N)^2.
\]
The finite matrix lower-frame theorem gives
\[
R_N^*G_{\mathcal X,N}R_N
\succeq
\gamma_Nc_{{\rm hyb},N}G_{R,N}
\succeq
\gamma_N\bigl(1-E_N(X_N)^2\bigr)G_{R,N}.
\]
If the last coefficient diverges and the finite residual windows promote through a fixed/exhaustive tail record, then the residual anti-invariant currency collapses on the completed residual sector.  This is precisely the source-absorption record for \(\Xi^{\rm BC}\). \end{proof}

\begin{corollary}[Maximal admissible shadow]
The best available choice under the source cap is \(X_N=Q_N^\theta\).  Thus a sufficient condition is
\[
\epsilon_{B,N}+\epsilon_{\rm Mell,N}+\kappa_{\rm tail,N}\to0,
\]
\[
A_N(1+T_N)^\alpha Q_N^{-\theta\eta}\to0,
\]
and
\[
\gamma_N\to\infty.
\]
Then
\[
c_{{\rm hyb},N}\to1,
\qquad
\gamma_Nc_{{\rm hyb},N}\to\infty.
\]
\end{corollary}

\begin{corollary}[Polynomial-growth phase boundary]
Suppose
\[
A_N\asymp Q_N^a,
\qquad
T_N\asymp Q_N^t.
\]
Then the Müntz contour residual can vanish under the source length cap only if
\[
\boxed{a+\alpha t<\theta\eta.}
\]
The boundary case \(a+\alpha t=\theta\eta\) gives no vanishing without additional logarithmic savings; the supercritical case
\[
a+\alpha t>\theta\eta
\]
forces a persistent readability defect under the declared source cap.
\end{corollary}

\begin{corollary}[Positive visibility floor]
Even if the visibility defect does not vanish, the source route may still close if there is a quantitative floor
\[
E_N(X_N)\le 1-\sigma_N,
\]
with
\[
\gamma_N\sigma_N\to\infty.
\]
Indeed
\[
1-E_N^2\ge 1-(1-\sigma_N)^2\ge \sigma_N.
\]
If \(\sigma_N\) decays too quickly for \(\gamma_N\) to compensate, the route remains support-only.
\end{corollary}

\section{Interpretation}
Step 123 turns the route into a three-way quantitative obligation:
\[
\boxed{\text{Burnol side: }\epsilon_{B,N}\to0,}
\]
\[
\boxed{\text{shadow side: }A_N(1+T_N)^\alpha X_N^{-\eta}\text{ small,}}
\]
\[
\boxed{\text{source side: }X_N\le Q_N^\theta\text{ and }\gamma_N\text{ grows}.}
\]
The decisive inequality is
\[
\boxed{A_N(1+T_N)^\alpha\ll Q_N^{\theta\eta}.}
\]
It says that the Burnol/Müntz window cannot be allowed to grow faster than the admissible short-polynomial length supplied by the source theorem.

\section{Nonclaims}
This step does not prove RH.  It does not prove the Heap--Soundararajan matrix lift.  It does not prove Burnol atom exhaustion.  It does not prove the Müntz angular gap.  It proves the conditional quantitative compatibility theorem that these ingredients must satisfy.

\end{document}
